% Source: https://artofproblemsolving.com/community/c3h2795209_2021_tx_mmc__texas_mathworks_math_contest
% Stable identity: aops-topic-2795209-q09
Write $n=2^k m$ with $m$ odd. If $k=0$ there are no even divisors, and $0$ is even. For $k\ge1$ an even divisor is $2^j d$ with $1\le j\le k$ and $d\mid m$, so there are $k\,\tau(m)$ of them, where $\tau$ counts divisors. This is odd exactly when $k$ is odd and $\tau(m)$ is odd, i.e. $m$ is a perfect square.

Then $n=2\cdot 2^{k-1}m=2\left(2^{\dfrac{k-1}{2}}\sqrt m\right)^2$, so $n$ is twice a perfect square; conversely every $n=2t^2$ has this form. Between $100$ and $300$ these are $2t^2$ for $t=8,\dots,12$:
\[128+162+200+242+288=\boxed{1020}.\]
